%% ============================================================
%%  chapitre4.tex — الفصل الرابع: التزامات شركة التشغيل المؤقت
%%                  والمقاولة المستعملة
%% ============================================================

\chapter{التزامات شركة التشغيل المؤقت والمقاولة المستعملة}
\label{chap:chapitre4}

\section*{مقدمة الفصل}

تقوم علاقة الشغل المؤقت على بنية ثلاثية تتوزع فيها صفة المشغل بين طرفين: فمقاولة التشغيل المؤقت هي المشغل القانوني الذي يبرم العقد ويؤدي الأجر، بينما تمارس المقاولة المستعملة سلطة الإدارة والتوجيه اليومي داخل منشأتها. ولتفادي أن يضيع الأجير بين الطرفين، وزّع المشرع المغربي الالتزامات بينهما توزيعاً دقيقاً في المواد 477 إلى 506 من مدونة الشغل.

ويجيب هذا الفصل عن السؤال الرابع من أسئلة النقاش: \textbf{ما الالتزامات القانونية لشركة التشغيل المؤقت تجاه الأجراء؟} فنعرض التزامات مقاولة التشغيل المؤقت (أولاً)، ثم التزامات المقاولة المستعملة (ثانياً)، ثم الضوابط الشكلية للعقود ومدتها وتجديدها (ثالثاً)، لنختم بجرد حقوق الأجير المؤقت وضماناته (رابعاً). وتفادياً للتكرار، نحيل على الفصل الثاني بخصوص التفاصيل التي سبق عرضها.

%% ─────────────────────────────────────────────────────────
\section{التزامات شركة التشغيل المؤقت}

تعرّف المادة 495 مقاولة التشغيل المؤقت بأنها شخص معنوي مستقل عن السلطة العمومية، يقتصر نشاطه على وضع أجراء رهن إشارة الغير بصفة مؤقتة، وهي التي \textbf{تشغّل هؤلاء الأجراء وتلتزم بأداء أجورهم وبالوفاء بالالتزامات القانونية الناشئة عن عقد الشغل}. ويترتب عن ذلك أن صفة المشغل القانوني تبقى لها طيلة مدة المهمة، ولو كان الأجير يعمل فعلياً داخل منشأة أخرى.

\subsection{الالتزامات السابقة لمزاولة النشاط}
لا يجوز لأي شركة أن تضع أجراء رهن إشارة الغير إلا إذا استوفت شروطاً مسبقة، سبق تفصيلها في الفصل الثاني، ونذكّر بأهمها:
\begin{itemize}[leftmargin=*, itemsep=4pt]
    \item \textbf{الترخيص المسبق} من السلطة الحكومية المكلفة بالشغل (المادة 477).
    \item \textbf{رأسمال لا يقل عن 100.000 درهم}، وانتفاء الأحكام الماسة بالشرف أو بالأمانة في حق المسيرين (المادة 481).
    \item \textbf{إيداع كفالة} لدى صندوق الإيداع والتدبير تعادل خمسين مرة القيمة الإجمالية السنوية للحد الأدنى القانوني للأجر (المادة 482).
    \item \textbf{حصر النشاط} في التشغيل المؤقت دون غيره (المادة 495).
\end{itemize}

\subsection{الالتزامات تجاه الأجير المؤقت}
\begin{enumerate}[leftmargin=*, itemsep=6pt]
    \item \textbf{أداء الأجر والوفاء بالالتزامات القانونية للمشغل} (المادة 495)، بما فيها التصريح بالأجير لدى الصندوق الوطني للضمان الاجتماعي (\textfrench{CNSS}).
    \item \textbf{إبرام عقد مكتوب مع الأجير} يتضمن البيانات الإلزامية التي تحددها المادة 501 (انظر القسم الثالث من هذا الفصل).
    \item \textbf{احترام المدد القصوى لفترة التجربة}: يومان إذا كان العقد أقل من شهر، وثلاثة أيام إذا تراوحت مدته بين شهر وشهرين، وخمسة أيام إذا تجاوز شهرين (المادة 502).
    \item \textbf{مجانية الخدمة}: يُمنع على الوكالة أن تتقاضى من الباحث عن الشغل أي أجر أو أتعاب، بصفة مباشرة أو غير مباشرة، كلياً أو جزئياً (المادة 480). كما يُمنع على مسيريها تسلّم أي ودائع أو ضمانات بمناسبة التشغيل (المادة 485).
    \item \textbf{عدم التمييز}: يُحظر كل تمييز مبني على العرق أو اللون أو الجنس أو الدين أو الرأي السياسي أو الأصل الوطني أو الاجتماعي أو الانتماء النقابي (المادة 478).
    \item \textbf{حماية المعطيات الشخصية}: لا تعالج الوكالة من معطيات الباحثين عن الشغل إلا ما يتعلق بمؤهلاتهم وخبرتهم المهنية، مع احترام حياتهم الخاصة (المادة 479).
\end{enumerate}

\subsection{الالتزامات تجاه الإدارة}
\begin{itemize}[leftmargin=*, itemsep=4pt]
    \item توجيه بيان مفصّل عن خدماتها إلى المصالح الإقليمية للتشغيل عند نهاية كل نصف سنة (المادة 484).
    \item مسك سجل وفق النموذج الذي تحدده السلطة المكلفة بالشغل، يُمكّن مفتش الشغل من مراقبة احترام القانون (المادة 486).
    \item بقاء التزاماتها تجاه الأجراء والصندوق الوطني للضمان الاجتماعي قائمة حتى بعد سحب الترخيص (المادة 503).
\end{itemize}

%% ─────────────────────────────────────────────────────────
\section{التزامات المقاولة المستعملة}

لا يعفي وجود وسيط مرخّص المقاولةَ المستعملة من كل التزام؛ فبما أنها هي التي تستفيد من العمل وتمارس سلطة التوجيه، أخضعها المشرع لالتزامات تسبق المهمة وأخرى تواكبها.

\subsection{الالتزامات السابقة للجوء إلى الأجير المؤقت}
\begin{enumerate}[leftmargin=*, itemsep=6pt]
    \item \textbf{استشارة الهيئات التمثيلية للأجراء} داخل المقاولة قبل اللجوء إلى التشغيل المؤقت (المادة 496).
    \item \textbf{التقيد بالحالات الحصرية الأربع} المنصوص عليها في المادة 496 (موضوع الفصل الثالث)، وعدم استعمال الأجير المؤقت لسدّ حاجيات دائمة.
    \item \textbf{عدم إسناد أشغال تكتسي خطورة خاصة} إلى الأجراء المؤقتين (المادة 497).
    \item \textbf{عدم اللجوء إلى الأجراء المؤقتين لمواجهة تزايد مؤقت في النشاط} داخل المناصب التي شملها فصل لأسباب اقتصادية، وذلك خلال السنة الموالية للفصل (المادة 498).
    \item \textbf{إبرام عقد مكتوب مع مقاولة التشغيل المؤقت} يبيّن سبب اللجوء ومدة المهمة ومكان تنفيذها ومقابل وضع الأجير رهن الإشارة (المادة 499).
\end{enumerate}

\subsection{الالتزامات أثناء تنفيذ المهمة}
\begin{itemize}[leftmargin=*, itemsep=6pt]
    \item \textbf{الصحة والسلامة}: اتخاذ جميع تدابير الوقاية والحماية لفائدة الأجراء المؤقتين، والتأمين عليهم ضد حوادث الشغل والأمراض المهنية (المادة 504). وهذا هو الالتزام الوحيد الذي حمّله المشرع صراحةً للمقاولة المستعملة بصفة مباشرة.
    \item \textbf{احترام المدة القصوى للمهمة} المحددة في المادة 500، وعدم الإبقاء على الأجير بعد انقضائها تحت غطاء تجديدات متتالية.
    \item \textbf{تطبيق قواعد تنظيم العمل السارية في المنشأة} (ساعات العمل، الراحة، حفظ الصحة)، بحكم أن الأجير يشتغل تحت إدارتها اليومية.
\end{itemize}

\begin{tcolorbox}[colback=lightgray, colframe=maincolor, arc=4pt,
    title={\textbf{تنبيه: المساواة في الأجر}}]
لا تتضمن مدونة الشغل نصاً صريحاً يُلزم المقاولة المستعملة بأن تمنح الأجير المؤقت أجراً مساوياً لأجر الأجير الدائم الذي يشغل المنصب نفسه، خلافاً للتشريع الفرنسي (المادة \textfrench{L1251-18} من قانون الشغل الفرنسي). أما المادة 9 فتحظر التمييز في الأجر على أسس محددة حصراً (العرق، اللون، الجنس، الإعاقة، الحالة الزوجية، الدين، الرأي السياسي، الانتماء النقابي، الأصل الوطني أو الاجتماعي)، وليس من بينها طبيعة العقد. لذلك يبقى الفارق في الأجر، كما في وضعية المجموعة 8، مسألةً يُعالجها الاجتهاد القضائي (الفصل التاسع) أو إعادة تكييف العلاقة (الفصل الخامس).
\end{tcolorbox}

%% ─────────────────────────────────────────────────────────
\section{العقود الكتابية ومدتها وتجديدها}

يقوم الشغل المؤقت على عقدين مكتوبين مترابطين: عقد وضع الأجير رهن الإشارة بين المقاولتين، وعقد المهمة بين مقاولة التشغيل المؤقت والأجير. ويوضح الجدول الآتي البيانات الإلزامية لكل منهما.

\begin{table}[h]
\centering
\small
\renewcommand{\arraystretch}{1.5}
\begin{tabular}{|p{3.2cm}|p{5.2cm}|p{5.2cm}|}
\hline
 & \textbf{عقد وضع الأجير رهن الإشارة (المادة 499)} & \textbf{عقد المهمة مع الأجير (المادة 501)} \\
\hline
\textbf{الأطراف} & مقاولة التشغيل المؤقت والمقاولة المستعملة & مقاولة التشغيل المؤقت والأجير \\
\hline
\textbf{الشكل} & مكتوب وجوباً & مكتوب وجوباً \\
\hline
\textbf{البيانات الإلزامية} & سبب اللجوء إلى الأجير المؤقت؛ مدة المهمة؛ مكان تنفيذها؛ مقابل وضع الأجير رهن الإشارة & بيانات المادة 499؛ المؤهلات المطلوبة؛ الأجر وكيفية أدائه؛ فترة التجربة؛ خصائص المنصب؛ رقم انخراط المقاولة ورقم تسجيل الأجير في الضمان الاجتماعي؛ شرط الترحيل إذا نُفّذت المهمة خارج المغرب \\
\hline
\textbf{بند خاص} & -- & إمكانية تشغيل الأجير من طرف المقاولة المستعملة بعد انتهاء المهمة \\
\hline
\end{tabular}
\caption{البيانات الإلزامية لعقدي الشغل المؤقت}
\label{tab:contrats-temporaire}
\end{table}

\subsection{المدة والتجديد}
ربطت المادة 500 مدة المهمة بسبب اللجوء إليها (انظر الجدول \ref{tab:durees-temporaire} في الفصل الثاني):
\begin{itemize}[leftmargin=*, itemsep=4pt]
    \item \textbf{إحلال أجير محل آخر}: لا تتجاوز المهمة مدة توقف عقد الأجير المعوَّض.
    \item \textbf{تزايد مؤقت في النشاط}: ثلاثة أشهر قابلة للتجديد \textbf{مرة واحدة فقط}، أي ستة أشهر كحد أقصى.
    \item \textbf{الأشغال الموسمية والأشغال التي جرى العرف بعدم إبرام عقود غير محددة المدة بشأنها}: ستة أشهر \textbf{غير قابلة للتجديد}.
\end{itemize}

\subsection{جزاء تجاوز المدة}
لم ينص المشرع صراحة، في الباب المتعلق بمقاولات التشغيل المؤقت، على تحويل عقد المهمة إلى عقد غير محدد المدة عند تجاوز هذه الآجال. ويترتب عن هذا السكوت أن الجزاء المدني يُستمدّ من الاجتهاد القضائي الذي يعتدّ بالواقع الفعلي للعلاقة، وهو ما نتناوله في الفصلين الخامس والتاسع، بينما يبقى الجزاء الزجري محصوراً في غرامة المادة 505.

%% ─────────────────────────────────────────────────────────
\section{حقوق الأجير المؤقت وضماناته}

تترجم الالتزامات السابقة إلى حقوق يمكن للأجير المؤقت أن يتمسك بها في مواجهة كل طرف حسب موقعه:

\begin{table}[h]
\centering
\small
\renewcommand{\arraystretch}{1.5}
\begin{tabular}{|p{4.2cm}|p{5.6cm}|p{3.6cm}|}
\hline
\textbf{الحق} & \textbf{مضمونه} & \textbf{الملزَم به والسند} \\
\hline
الأجر & أجر لا يقل عن الحد الأدنى القانوني، يؤدى في آجاله & مقاولة التشغيل المؤقت (495) \\
\hline
الحماية الاجتماعية & التصريح الكامل بالأجير لدى \textfrench{CNSS} عن كل أيام العمل & مقاولة التشغيل المؤقت (495، 501) \\
\hline
الصحة والسلامة & تدابير الوقاية والتأمين ضد حوادث الشغل والأمراض المهنية & المقاولة المستعملة (504) \\
\hline
العقد المكتوب & عقد يحدد المهمة والأجر والمدة & مقاولة التشغيل المؤقت (501) \\
\hline
المجانية & عدم أداء أي مقابل للوكالة & مقاولة التشغيل المؤقت (480، 485) \\
\hline
عدم التمييز & المساواة في الولوج إلى الشغل & مقاولة التشغيل المؤقت (478) \\
\hline
ضمان الديون & استيفاء الأجور والاشتراكات من الكفالة عند عجز الوكالة & صندوق الإيداع والتدبير (482، 488) \\
\hline
التشغيل الدائم & إمكانية التشغيل المباشر لدى المقاولة المستعملة بعد المهمة & بند إلزامي في العقد (501) \\
\hline
\end{tabular}
\caption{حقوق الأجير المؤقت والجهة الملزمة بها}
\label{tab:droits-interimaire}
\end{table}

ويُضاف إلى ذلك أن تقديم الأجير لشكاية ضد مشغله أمام مفتش الشغل لا يُعدّ مبرراً مقبولاً لفصله (المادة 36)، وهي ضمانة أساسية تتيح للأجير المؤقت الدفاع عن حقوقه دون خوف من فقدان عمله.

\subsection{إسقاط على وضعية المجموعة 8}
بإسقاط ما سبق على الوضعية المدروسة، يتبيّن أن العمال الذين توفرهم شركة المناولة يشتغلون في المصنع نفسه \textbf{منذ سنوات}، ويؤدّون المهام ذاتها التي يقوم بها أجراء الشركة الأصلية، مقابل أجور أقل وامتيازات اجتماعية محدودة. فإذا كُيّفت العلاقة على أنها شغل مؤقت، فإن المقاولتين معاً قد أخلّتا بالتزاماتهما: تجاوز المدة القصوى للمادة 500 بشكل صارخ، وتلبية حاجيات دائمة خارج الحالات الحصرية للمادة 496. أما إن قُدّمت على أنها مناولة، فإن اقتصار موضوع العقد على «توفير اليد العاملة» يجعلها نشاطاً لا يحق إلا لمقاولة تشغيل مؤقت مرخّصة. وفي الحالتين، يبقى السؤال المطروح: \textbf{من يتحمل المسؤولية عن هذه الإخلالات؟} وهو موضوع الفصل الخامس.

\section*{خلاصة الفصل}
وزّع المشرع المغربي الالتزامات في الشغل المؤقت بين طرفين: مقاولة التشغيل المؤقت، بصفتها المشغل القانوني، تتحمل الأجر والتصريح الاجتماعي والعقد المكتوب والمجانية، وتضمن ديونها كفالةٌ مودعة لدى صندوق الإيداع والتدبير؛ والمقاولة المستعملة تتحمل احترام حالات اللجوء ومدده، والسلامة والتأمين ضد حوادث الشغل. غير أن هذا البناء يشكو من ثغرتين أساسيتين: غياب نص صريح يقرر المساواة في الأجر، وغياب جزاء مدني تلقائي عند تجاوز المدد القانونية. ولذلك ننتقل في الفصل الخامس إلى تحديد المسؤولية القانونية لكل من الشركة الأصلية وشركة المناولة عند الإخلال بهذه الالتزامات.
